
% LaTeX Beamer file automatically generated from DocOnce
% https://github.com/doconce/doconce

%-------------------- begin beamer-specific preamble ----------------------
\documentclass{beamer}

\usepackage{amsmath,amsfonts,amssymb}
\usepackage{physics}
\usepackage{bm}
\usepackage{quantikz}
\usepackage{listings}


\usetheme{red_plain}
\usecolortheme{default}

% turn off the almost invisible, yet disturbing, navigation symbols:
\setbeamertemplate{navigation symbols}{}

% Examples on customization:
%\usecolortheme[named=RawSienna]{structure}
%\usetheme[height=7mm]{Rochester}
%\setbeamerfont{frametitle}{family=\rmfamily,shape=\itshape}
%\setbeamertemplate{items}[ball]
%\setbeamertemplate{blocks}[rounded][shadow=true]
%\useoutertheme{infolines}
%
%\usefonttheme{}
%\useinntertheme{}
%
%\setbeameroption{show notes}
%\setbeameroption{show notes on second screen=right}

% fine for B/W printing:
%\usecolortheme{seahorse}

\usepackage{pgf}
\usepackage{graphicx}
\usepackage{epsfig}
\usepackage{relsize}

\usepackage{fancybox}  % make sure fancybox is loaded before fancyvrb

\usepackage{fancyvrb}
\usepackage{minted} % requires pygments and latex -shell-escape filename
%\usepackage{anslistings}
%\usepackage{listingsutf8}

\usepackage{amsmath,amssymb,bm}
%\usepackage[latin1]{inputenc}
\usepackage[T1]{fontenc}
\usepackage[utf8]{inputenc}
\usepackage{colortbl}
\usepackage[english]{babel}
\usepackage{tikz}
\usepackage{framed}
% Use some nice templates
\beamertemplatetransparentcovereddynamic

% --- begin table of contents based on sections ---
% Delete this, if you do not want the table of contents to pop up at
% the beginning of each section:
% (Only section headings can enter the table of contents in Beamer
% slides generated from DocOnce source, while subsections are used
% for the title in ordinary slides.)
\AtBeginSection[]
{
  \begin{frame}<beamer>[plain]
  \frametitle{}
  %\frametitle{Outline}
  \tableofcontents[currentsection]
  \end{frame}
}
% --- end table of contents based on sections ---

% If you wish to uncover everything in a step-wise fashion, uncomment
% the following command:

%\beamerdefaultoverlayspecification{<+->}

\newcommand{\shortinlinecomment}[3]{\note{\textbf{#1}: #2}}
\newcommand{\longinlinecomment}[3]{\shortinlinecomment{#1}{#2}{#3}}

\definecolor{linkcolor}{rgb}{0,0,0.4}
\hypersetup{
    colorlinks=true,
    linkcolor=linkcolor,
    urlcolor=linkcolor,
    pdfmenubar=true,
    pdftoolbar=true,
    bookmarksdepth=3
    }
\setlength{\parskip}{0pt}  % {1em}

\newenvironment{doconceexercise}{}{}
\newcounter{doconceexercisecounter}
\newenvironment{doconce:movie}{}{}
\newcounter{doconce:movie:counter}

\newcommand{\subex}[1]{\noindent\textbf{#1}}  % for subexercises: a), b), etc

\logo{{\tiny \copyright\ 1999-2026, Morten Hjorth-Jensen. Released under CC Attribution-NonCommercial 4.0 license}}

%-------------------- end beamer-specific preamble ----------------------

% Add user's preamble




% insert custom LaTeX commands...

\raggedbottom
\makeindex

%-------------------- end preamble ----------------------

\begin{document}

% matching end for #ifdef PREAMBLE

\newcommand{\exercisesection}[1]{\subsection*{#1}}



% ------------------- main content ----------------------



% ----------------- title -------------------------

\title{Quantum Computing and Quantum Machine Learning}

% ----------------- author(s) -------------------------

\author{Morten Hjorth-Jensen\inst{1}}
\institute{Department of Physics, University of Oslo, Norway\inst{1}}
% ----------------- end author(s) -------------------------

\date{March 18, 2026
% <optional titlepage figure>
\ \\ 
{\tiny \copyright\ 1999-2026, Morten Hjorth-Jensen. Released under CC Attribution-NonCommercial 4.0 license}
}

\begin{frame}[plain,fragile]
\titlepage
\end{frame}

\begin{frame}[plain,fragile]
\frametitle{Plans for the week of March 16-20, 2026}

\begin{block}{}
\begin{enumerate}
\item From classical Fourier transforms to quantum Fourier transforms (QFTs).

\item Reading recommendation Hundt, Quantum Computing for Programmers, sections 6.1-6.4 on QFT.

\item Discussion and work on  project 1
\end{enumerate}

\noindent
\end{block}
\end{frame}

\begin{frame}[plain,fragile]
\frametitle{Overarching motivation}

We will start with a reminder on Fourier theory and in particular on discrete Fourier transforms (DFTs).  There are many
motivations for the DFTs. For those of you familiar with signal
processing, harmonic oscillations, and many other areas of
applications, Fourier transforms are almost standard kitchen
items. 

But first some general motivation for Quantum Fourier Transforms (QFTs).
\end{frame}

\begin{frame}[plain,fragile]
\frametitle{Quantum Fourier Transforms (QFTs)}

\begin{itemize}
\item QFTs are the quantum analogue of the Discrete Fourier Transforms (DFTs).

\item They  play a crucial role in quantum algorithms like Shor’s Algorithm and Quantum Phase Estimation (algo for finding for example eigenvalues)

\item QFTs provide an \emph{exponential speedup} over classical Fourier Transform methods.
\end{itemize}

\noindent
\end{frame}

\begin{frame}[plain,fragile]
\frametitle{Why Quantum Fourier Transforms and exponential speedup}

\begin{itemize}
\item Classical Discrete Fourier Transforms (DFTs) require $O(N^2)$ operations.

\item The Fast Fourier Transform (FFT) improves this to $O(N \log N)$ operations.

\item QFTs reduce complexity to $O((\log N)^2)$ using quantum gates.
\end{itemize}

\noindent
\end{frame}

\begin{frame}[plain,fragile]
\frametitle{Quantum Fourier Transforms and quantum parallelism}

\begin{itemize}
\item QFTs act on a \emph{superposition} of states, processing all inputs simultaneously.

\item This is crucial in algorithms like:
\begin{itemize}

  \item Shor’s Algorithm for factoring large numbers.

  \item Quantum Phase Estimation (QPE) for eigenvalue extraction.
\end{itemize}

\noindent
\end{itemize}

\noindent
\end{frame}

\begin{frame}[plain,fragile]
\frametitle{Quantum Fourier Transforms and implementations}

\begin{itemize}
\item QFTs require only \emph{Hadamard gates and controlled-phase gates}.

\item A 3-qubit QFT circuit uses only $O(n^2)$.

\item This makes QFTs highly efficient for \emph{quantum hardware}.
\end{itemize}

\noindent
\end{frame}

\begin{frame}[plain,fragile]
\frametitle{Quantum Fourier Transform and Applications in Quantum Computing}

\begin{itemize}
\item \textbf{Shor’s Algorithm:} Uses QFTs to find periodicity in modular exponentiation.

\item \textbf{Quantum Phase Estimation (QPE):} QFTs extract eigenvalues of for example unitary matrices or finding eigenvalues of a given Hamiltonian.

\item \textbf{Quantum Signal Processing:} QFTs enable spectral analysis on quantum data. Important for AI applications.
\end{itemize}

\noindent
\end{frame}

\begin{frame}[plain,fragile]
\frametitle{Quantum Fourier Transforms and Reduced Memory Complexity (not covered)}

\begin{itemize}
\item Classical DFTs require $O(N)$ space.

\item QFT store Fourier-transformed coefficients \emph{implicitly} in qubit amplitudes.

\item Only $O(n)$ qubits are needed for a size $N = 2^n$ transformation.
\end{itemize}

\noindent
\end{frame}

\begin{frame}[plain,fragile]
\frametitle{Quantum Fourier Transforms and Quantum Signal Processing (not covered in this course)}

QFTs enable applications such as:
\begin{itemize}
\item Quantum spectral analysis.

\item Quantum image processing.

\item Quantum filtering and denoising.
\end{itemize}

\noindent
These can enhance AI, cryptography, and data processing.
\end{frame}

\begin{frame}[plain,fragile]
\frametitle{Why Quantum Fourier Transforms? Brief summary}

\begin{itemize}
\item Quantum Fourier Transforms provide \emph{exponential speedup}

\item They enable key quantum algorithms like \emph{Shor’s Algorithm} and \emph{Quantum Phase Estimation}.

\item QFTs are  more efficient in \emph{memory usage and circuit complexity} than classical methods.

\item Future quantum applications in \emph{signal processing, AI, and cryptography} will heavily rely on QFT.
\end{itemize}

\noindent
\end{frame}

\begin{frame}[plain,fragile]
\frametitle{Now technicalities: Reminder on Fourier theory, a familiar case first}

For problems with so-called harmonic oscillations, given by for example the following differential equation
\[
m\frac{d^2x}{dt^2}+\eta\frac{dx}{dt}+x(t)=f(t),
\]
where $f(t)$ is an applied external force acting on the system (often
called a driving force), one can use the theory of Fourier
transformations to find the solutions of this type of equations.
\end{frame}

\begin{frame}[plain,fragile]
\frametitle{Several driving forces}

If one has several driving forces, $f(t)=\sum_n f_n(t)$, one can find
the particular solution $x_{pn}(t)$ to the above differential equation for each $f_n$. The particular
solution for the entire driving force is then given by a series like

\[
x_p(t)=\sum_nx_{pn}(t).
\]

This is known as the principle of superposition. It only applies when
the homogenous equation is linear. 
Superposition is especially useful when $f(t)$ can be written
as a sum of sinusoidal terms, because the solutions for each
sinusoidal (sine or cosine)  term is analytic.
\end{frame}

\begin{frame}[plain,fragile]
\frametitle{Periodicity}

Driving forces are often periodic, even when they are not
sinusoidal. Periodicity implies that for some time $t$ our function repeats itself periodically after a period $T$, that is

\[
f(t+\tau)=f(t). 
\]

One example of a non-sinusoidal periodic force is a square wave. Many
components in electric circuits are non-linear, for example diodes. This 
makes many wave forms non-sinusoidal even when the circuits are being
driven by purely sinusoidal sources.
\end{frame}

\begin{frame}[plain,fragile]
\frametitle{Continuous Fourier transforms and the principle of Superposition}

It was Fourier's idea to expand a continuous and periodic function $f(t)$ in terms of sums sine and cosine ordered functions (we will use exponentials however)
as
\[
f(t) = \sum_{k=1}^{n}a_k\sin{(2\pi kt+\phi_n)},
\]
with $\phi_n$ being a constant phase. The function $f$ is assumed to be bounded in order to be able to define properly the error in truncating the sum
over $n$, that is
\[
\int_a^b dt \vert f(t)\vert^2 \le M < \infty.
\]

Below we discuss how to find the coefficients $a_n$.
\end{frame}

\begin{frame}[plain,fragile]
\frametitle{Rewriting in terms of sines and cosines}

It is common to rewrite the above sum in terms of sines and cosines and to add a complex constant $a_0$. This gives us
\[
f(t) = \frac{a_0}{2}+ \sum_{k=1}^{n}\left(a_k\cos{(2\pi kt)}+b_n\sin{(2\pi kt)}\right).
\]
Using the standard trigonometric relations
\begin{align*}
\cos{t} = \frac{\exp{\imath t}+\exp{-\imath t}}{2} & \hspace{0.5cm} \sin{t} = \frac{\exp{\imath t}-\exp{-\imath t}}{2\imath}.
\end{align*}
\end{frame}

\begin{frame}[plain,fragile]
\frametitle{Exponential expression}

We can rewrite the Fourier expansion as
\[
f(t) =\sum_{k=-n}^{n}c_{k}\exp{(2\pi\imath kt)}, 
\]
with $c_0=a_0/2$. The coefficients $c_k$ are complex and satisfy $c_{-k}=c_k^*$.
The constant $c_{0}=c_0^*$ is a real number.

The above sum can also be rewritten in terms of the real part of the exponentials as
\[
f(t) =2\mathrm{Re}\left(\sum_{k=0}^{n}c_{k}\exp{(2\pi\imath kt} \right),
\]
where we used that $c_{-k}=c_k^*$. We leave it as an exercise to the reader to show the latter expression.

How do we determine the coefficients $c_k$?
\end{frame}

\begin{frame}[plain,fragile]
\frametitle{Determining $c_k$}

Let us assume that we have a periodic function $f(t)$ 
\[
f(t) =\sum_{k=-n}^{n}c_{k}\exp{(2\pi\imath kt)}, 
\]
and we select the coefficient $c_l$, and multiply the l.h.s and r.h.s. with $\exp{(-2\pi\imath lt)}$, and isolate the $c_l$ terms
\[
c_l=f(t)\exp{(-2\pi\imath lt)}-\sum_{k=-n, k\ne l}^{n}c_{k}\exp{(2\pi\imath (k-l)t)}. 
\]
\end{frame}

\begin{frame}[plain,fragile]
\frametitle{Integrating both sides}

Then we integrate both sides (note we have assumed a period of one) from $0$ to $1$. This gives
\[
c_l=\int_0^1 dtf(t)\exp{(-2\pi\imath lt)},
\]
since the terms with the sum is zero.  We leave again this derivation as an exercise to the dedicated reader.
If the function $f$ is real ($t$ nd $dt$ are real), then $c_l^*=c_{-l}$.
The coefficients $c_n$ are normally called the Fourier coefficients and it is common to relabel them in terms of the function $f$ as
\[
\hat{f}(k)=\int_0^1 dt f(t)\exp{(-2\pi\imath kt)}.
\]
\end{frame}

\begin{frame}[plain,fragile]
\frametitle{Independence of interval length}

In the discussions and equations above, we have assumed that the
integration interval has a length of one, that is a period of length one. It is easy to change this
length and as we show below, if we integrate from say $a$ to $a+1$, this is the same as integrating over an interval of length $1$.

To show this, we need also to recall that our function $f(t)$ is assumed to be periodic, that is we need to satisfy
\[
f(a+1)=f(a).
\]
Later we will generalize this to a period of arbitrary length.

Assume that $a$ is any number and taking the derivative of the integral with respect to $a$, we have
\[
\frac{d}{da}\left[\int_a^{a+1} dt f(t)\exp{(-2\pi\imath kt)}\right]. 
\]
\end{frame}

\begin{frame}[plain,fragile]
\frametitle{Writing out the various terms}

From the last equation we have then
\[
\exp{(-2\pi\imath ka)}\exp{(-2\pi\imath k)}f(a+1)-\exp{(-2\pi\imath ka)}f(a)=0,
\]
where we used the periodicity  $f(a+1)=f(a)$ and that $\exp{(-2\pi\imath k)}=1$ since $k$ is an integer.
This shows that the expression for $\hat{f}$ is independent of $a$. A common instance is
\[
\hat{f}(k)=\int_{-\frac{1}{2}}^{\frac{1}{2}} dt f(t)\exp{(-2\pi\imath kt)}.
\]
\end{frame}

\begin{frame}[plain,fragile]
\frametitle{Changing period}

What if the period is not equal to one? Assume we are working with a
function $f(t)$ whose period is $T$. We can then define a function
$g(t)$ with period one as
\[
g(t) = f(Tt),
\]
that is 
\[
g(t) =\sum_{k=-n}^{n}c_{n}\exp{(2\pi\imath kt)}, 
\]
and introducing the variable $s=Tt$ we have $g(t)=f(s)$ we have
\[
f(s) = g(t)=\sum_{k=-n}^{n}c_{n}\exp{(2\pi\imath kt)}=\sum_{k=-n}^{n}c_{k}\exp{(2\pi\imath ks/T)}. 
\]
\end{frame}

\begin{frame}[plain,fragile]
\frametitle{Harmonics}

The so-called harmonics are now defined as $\exp{(2\pi\imath ks/T)}$.
We have 
\[
\hat{g}(k)=\int_0^1 dt g(t)\exp{(-2\pi\imath kt)},
\]
and using $s=Tt$ we have
\[
\hat{g}(k)=\int_0^1 dt g(t)\exp{(-2\pi\imath kt)}=\frac{1}{T}\int_0^T ds f(s)\exp{(-2\pi\imath ks/T)}.
\]

For a period $T$ we have thus the general expression for the Fourier transform
\[
\hat{f}(k)=c_k=\frac{1}{T}\int_0^T dt f(t)\exp{(-2\pi\imath kt/T)}.
\]
\end{frame}

\begin{frame}[plain,fragile]
\frametitle{Typical interval}

An often used choice of interval is
\[
\hat{f}(k)=c_k=\frac{1}{T}\int_{-T/2}^{T/2} dt f(t)\exp{(-2\pi\imath kt/T)},
\]
and a common choice for the harmonics is $(1/\sqrt{T})\exp{(2\pi\imath kt/T)}$.

As a small addendum, if the signal (function) $f(t)$ is real and even one has $f(-t)=f(t)$. If $f$ is an even function, then $\hat{f}$ is also an even function which leads to $\hat{f}(-k)=\hat{f}(k)$.
\end{frame}

\begin{frame}[plain,fragile]
\frametitle{Sinusoidal example}

For the sinusoidal example the
period is $T=2\pi/\omega$. However, higher harmonics can also
satisfy the periodicity requirement. In general, any force that
satisfies the periodicity requirement can be expressed as a sum over
harmonics,

\[
f(t)=\frac{a_0}{2}+\sum_{n>0} a_n\cos(2n\pi t/T)+b_n\sin(2n\pi t/T).
\]
\end{frame}

\begin{frame}[plain,fragile]
\frametitle{Square well example}

In the code discussed earlier, we used a square well as a our
example. If we assume our period has length $1$, we can define the
square well function $f(t)$ as

\[
f(t)=\left\{\begin{array}{cc} 1 & 0 \le t \le \frac{1}{2} \\
                              -1 & \frac{1}{2} < t \le 1\end{array}\right.
\]

The zeroth coefficient $a_0$ is zero since it is the average of the function $f(t)$ over the intergration domain $t\in [0,1]$.
\end{frame}

\begin{frame}[plain,fragile]
\frametitle{The coefficients}

We find that the other coefficients are
\[
\hat{f}(n)=c_n=\int_0^1 dt f(t)\exp{(-2\pi\imath nt)}=\frac{1}{\imath\pi n}\left(1-\exp{-\imath\pi n}\right).
\]

We note that $(1-\exp{-\imath\pi n})$ is zero when $n$ is an even
number and $2$ when $n$ is an odd number. We can then combine the
positive and negative values of $n$ using
\[
\exp{(2\pi\imath nt)}-\exp{(-2\pi\imath nt)}=2\imath \sin{(2\pi nt)},
\]
and we have
\[
f(t)=\frac{4}{\pi}\sum_{k=0}^{\infty}\frac{1}{2k+1}\sin{(2\pi(2k+1)t)}.
\]
\end{frame}


\begin{frame}[plain,fragile]
\frametitle{Inverse Fourier transform}

The inverse Fourier transform is given by
\[
\bm{F}^{-1}[g(y)]=\frac{1}{2\pi}\int_{-\infty}^{\infty} d\omega \exp{i\omega y} g(\omega).
\]

The inverse Fourier transform of the product of the two functions $\hat{f}\hat{g}$ can be written as
\[
\bm{F}^{-1}[(\hat{f}\hat{g})(x)]=\frac{1}{2\pi}\int_{-\infty}^{\infty} d\omega \exp{i\omega x} \hat{f}(\omega)\hat{g}(\omega).
\]
\end{frame}

\begin{frame}[plain,fragile]
\frametitle{Rewriting}

We can rewrite the latter as
\[
\bm{F}^{-1}[(\hat{f}\hat{g})(x)]=\int_{-\infty}^{\infty} d\omega \exp{i\omega x} \hat{f}(\omega)\left[\frac{1}{2\pi}\int_{-\infty}^{\infty}g(y)dy \exp{-i\omega y}\right]=\frac{1}{2\pi}\int_{-\infty}^{\infty}dy g(y)\int_{-\infty}^{\infty} d\omega \hat{f}(\omega) \exp{i\omega(x- y)},
\]
which is simply 
\[
\bm{F}^{-1}[(\hat{f}\hat{g})(x)]=\int_{-\infty}^{\infty}dy g(y)f(x-y)=(f*g)(x),
\]
the convolution of the functions $f$ and $g$.
\end{frame}

\begin{frame}[plain,fragile]
\frametitle{Transforming to discrete variables}

In the Fourier transform $c_n$ is transformed from a dicrete variable
to a continuous one as $T \rightarrow \infty$. We then replace $c_n$ with
$f(k)dk$ and let $n/T \rightarrow k$, and the sum is changed to an
integral. This gives
\[
f(x) = \int_{-\infty}^{\infty}dkF(k) \exp{i(2\pi kx)}
\]
and
\[
F(k) = \int_{-\infty}^{\infty}dxf(x) \exp{-i(2\pi kx)}
\]
One way to interpret the Fourier transform is then as a transformation from one basis to another.
\end{frame}

\begin{frame}[plain,fragile]
\frametitle{Discrete Fourier transform}

We change notation by replacing $t$ with $x$.

Next we make another generalization by having a discrete function,
that is $f(x) \rightarrow f(x_k)$ with $x_k = k\Delta x$ for $k=0,\dots, N-1$. This leads to the sums
\[
f_x = \frac{1}{n} \sum_{k=0}^{n-1}F_k \exp{i(2\pi kx)/n},
\]
and
\[
F_k = \sum_{x=0}^{n-1}f_x \exp{-i(2\pi kx)/n}.
\]

Although we have used functions here, this could also be a set of
numbers.
\end{frame}

\begin{frame}[plain,fragile]
\frametitle{Simple example}

As an example we can have a set of complex numbers
$\{x_0,\dots,x_{n-1}\}$ with fixed length $n$, we can Fourier
transform this as
\[
y_k = \frac{1}{\sqrt{n}} \sum_{j=0}^{n-1} x_j \exp{i(2\pi jk)/n},
\]

leading to a new set of complex numbers $\{ y_0,\dots,y_{n-1}\}$.
\end{frame}

\begin{frame}[plain,fragile]
\frametitle{Discrete Fourier Transformations}

Consider two sets of complex numbers $x_k$ and $y_k$ with
$k=0,1,\dots,n-1$ entries. The discrete Fourier transform is defined
as
\[
y_k = \frac{1}{\sqrt{n}} \sum_{j=0}^{n-1} \exp{(\frac{2\pi\imath jk}{n})} x_j.
\]
As an example, assume $x_0=1$ and $x_1=1$. We can then use the above expression to find $y_0$ and $y_1$.

With the above formula we get then
\begin{align*}
y_0 &= \frac{1}{\sqrt{2}} \left( \exp{(\frac{2\pi\imath 0\times 1}{2})} \times 1+\exp{(\frac{2\pi\imath 0\times 1}{2})}\times 2\right)\\
& =\frac{1}{\sqrt{2}}(1+2)=\frac{3}{\sqrt{2}},
\end{align*}
and
\begin{align*}
y_1 &= \frac{1}{\sqrt{2}} \left( \exp{(\frac{2\pi\imath 0\times 1}{2})} \times 1+\exp{(\frac{2\pi\imath 1\times 1}{2})}\times 2\right)=\\
& =\frac{1}{\sqrt{2}}(1+2\exp{(\pi\imath)})=-\frac{1}{\sqrt{2}}.
\end{align*}
\end{frame}

\begin{frame}[plain,fragile]
\frametitle{More details on Discrete Fourier transforms}

Suppose that we have a vector $f$ of $n$ complex numbers, $f_{k}, k
\in\{0,1, \ldots, n-1\}$. Then the discrete Fourier transform (DFT) is
a map from these $n$ complex numbers to $n$ complex numbers, the
Fourier transformed coefficients $\tilde{f}_{j}$, given by

\begin{equation*}
\tilde{f}_{j}=\frac{1}{\sqrt{n}} \sum_{k=0}^{n-1} \omega^{-j k} f_{k}
\end{equation*}
where $\omega=\exp \left(\frac{2 \pi i}{N}\right)$.
\end{frame}

\begin{frame}[plain,fragile]
\frametitle{Inverse DFT}

The inverse DFT is given by

\begin{equation*}
f_{j}=\frac{1}{\sqrt{n}} \sum_{k=0}^{n-1} \omega^{j k} \tilde{f}_{k}
\end{equation*}

To see this consider how the basis vectors transform. If $f_{k}^{l}=\delta_{k, l}$, then

\begin{equation*}
\tilde{f}_{j}^{l}=\frac{1}{\sqrt{n}} \sum_{k=0}^{n-1} \omega^{-j k} \delta_{k, l}=\frac{1}{\sqrt{n}} \omega^{-j l}
\end{equation*}
\end{frame}

\begin{frame}[plain,fragile]
\frametitle{Orthonormality}

These DFT vectors are orthonormal, that is
\begin{equation*}
\sum_{j=0}^{n-1} \tilde{f}^{l}{ }_{j}^{*} \tilde{f}_{j}^{m}=\frac{1}{n} \sum_{j=0}^{n-1} \omega^{j l} \omega^{-j m}=\frac{1}{N} \sum_{j=0}^{n-1} \omega^{j(l-m)} 
\end{equation*}

This last sum can be evaluated as a geometric series, but beware of the $(l-m)=0$ term, and yields

\begin{equation*}
\sum_{j=0}^{n-1} \tilde{f}^{l}{ }_{j}^{*} \tilde{f}_{j}^{m}=\delta_{l, m}
\end{equation*}
\end{frame}

\begin{frame}[plain,fragile]
\frametitle{Inverse transform}

From this we can check that the inverse DFT does indeed perform the inverse transform:

\begin{equation*}
f_{j}=\frac{1}{\sqrt{n}} \sum_{k=0}^{n-1} \omega^{j k} \tilde{f}_{k}=\frac{1}{\sqrt{n}} \sum_{k=0}^{n-1} \omega^{j k} \frac{1}{\sqrt{n}} \sum_{l=0}^{n-1} \omega^{-l k} f_{l}=\frac{1}{n} \sum_{k, l=0}^{n-1} \omega^{(j-l) k} f_{l}=\sum_{l=0}^{n-1} \delta_{j, l} f_{l}=f_{j} 
\end{equation*}
\end{frame}

\begin{frame}[plain,fragile]
\frametitle{Summary: Discrete Fourier transforms}

The Discrete Fourier Transform is a tool used in many different fields
and finds applications in for example signal processing.  It has many
many applications (in Speech Recognition, Data Processing, Optics,
Acoustics). It is used in any application which performs Digital
Signal Processing.

There
is an efficient algorithm, called the Fast Fourier Transform (FFT),
which uses the special relationship amongst the roots of unity to
compute the entire DFT in $O(n \log{n})$ time. Moreover, it is possible to take the DFT and recover the original
polynomial in its coefficient form in $O(n \log{n})$ time, by using the
FFT to compute the inverse DFT.
\end{frame}

\begin{frame}[plain,fragile]
\frametitle{Fast Fourier transform (FFT)}

Because of the many important applications of DFT, the Fast Fourier
Transform is ranked among the top ten  algorithms in the 20th century.
It came to light as a tool for Signal processing in 1965, when
Cooley (of IBM) and Tukey (of Princeton) wrote a paper presenting the
algorithm. However the basics of the algorithm were known long before
that, for example, it was known to Gauss back in 1805.
To read more about Fast Fourier transforms and similar topics, see for example \href{{https://link.springer.com/book/10.1007/978-1-4020-6629-0}}{Fast Fourier Transform - Algorithms and Applications}. See also \href{{https://github.com/CompPhysics/QuantumComputingMachineLearning/blob/gh-pages/doc/Textbooks/fastfourier.pdf}}{\nolinkurl{https://github.com/CompPhysics/QuantumComputingMachineLearning/blob/gh-pages/doc/Textbooks/fastfourier.pdf}}

Our emphasis is on the link between discrete Fourier transforms and quantum Fourier transforms.
For fast Fourier transforms, we may (if there is interest) have a dedicated sessions.
\end{frame}

\begin{frame}[plain,fragile]
\frametitle{More summary: The discrete Fourier transform (DFT)}

The discrete Fourier transform takes as input a complex vector
\begin{align*}
\bm{x} =\vert\bm{x}\rangle= 
\begin{bmatrix}
x_0 \\ 
x_1 \\ 
x_2 \\
\vdots \\ 
x_{n-2} \\
x_{n-1}
\end{bmatrix}.
\end{align*}

Each entry of the output vector $\vert \bm{y}\rangle$ is given by

\begin{align*}
y_k = \frac{1}{\sqrt{N}} \sum_{j=0}^{N-1} x_je^{2 \pi ijk/N}.
\end{align*}
\end{frame}

\begin{frame}[plain,fragile]
\frametitle{Output vector}

The output is a complex vector

\begin{align*}
\vert \bm{y}\rangle = 
\frac{1}{\sqrt{n}}
\begin{bmatrix}
\sum_{j=0}^{n-1} x_je^{2\pi ij0/n} \\
\sum_{j=0}^{n-1} x_je^{2\pi ij1/n} \\
\sum_{j=0}^{n-1} x_je^{2\pi ij2/n} \\
\vdots \\
\sum_{j=0}^{n-1} x_je^{2\pi ij(n-2)/n} \\
\sum_{j=0}^{n-1} x_je^{2\pi ij(n-1)/n}
\end{bmatrix}
\end{align*}
\end{frame}

\begin{frame}{Summarizing Discrete Fourier Transform}

For vector

\[
x=(x_0,\dots,x_{N-1})
\]

\[
y_k =
\frac{1}{\sqrt N}
\sum_{j=0}^{N-1}
x_j e^{2\pi i jk/N}
\]

\end{frame}

%------------------------------------------------

\begin{frame}{Fourier Matrix}

\[
F_{jk}=
\frac{1}{\sqrt N}
e^{2\pi i jk/N}
\]

Properties

\begin{itemize}
\item unitary
\item periodic
\item diagonalizes cyclic shifts
\end{itemize}

\end{frame}



\begin{frame}[plain,fragile]
\frametitle{From DFT to QFT}

The Quantum Fourier Transform (QFT) has mathematically the same equation as starting point
but the notation is generally different. We generally compute the
quantum Fourier transform on a set of orthonormal basis state vectors
\(|0\rangle, |1\rangle, ..., |N-1\rangle\). The linear operator defining
the transform is given by the action on basis states

\[
|j\rangle \mapsto \frac{1}{\sqrt{N}} \sum_{k=0}^{N-1} e^{2\pi ijk/N}|k\rangle.
\]
\end{frame}

\begin{frame}[plain,fragile]
\frametitle{From DFT to QFT}

The Quantum Fourier Transform (QFT) has mathematically the same equation as starting point
but the notation is generally different. We generally compute the
quantum Fourier transform on a set of orthonormal basis state vectors
\(|0\rangle, |1\rangle, ..., |N-1\rangle\). The linear operator defining
the transform is given by the action on basis states

\[
|j\rangle \mapsto \frac{1}{\sqrt{N}} \sum_{k=0}^{N-1} e^{2\pi ijk/N}|k\rangle.
\]
\end{frame}

\begin{frame}[plain,fragile]
\frametitle{In terms of arbitrary states}

This can be written on arbitrary states,
\[
\sum_{j=0}^{N-1}x_j|j\rangle \mapsto \sum_{k=0}^{N-1} y_k|k\rangle
\]
where each amplitude \(y_k\) is the discrete Fourier transform of
\(x_j\).
\end{frame}

\begin{frame}[plain,fragile]
\frametitle{Unitarity}

The quantum Fourier transform is unitary. Taking \(N = 2^n\),
for \(n\) qubits gives us the orthonormal (computational) basis
\[
|0\rangle, |1\rangle, ..., |2^{n}-1\rangle.
\]
\end{frame}

\begin{frame}[plain,fragile]
\frametitle{Full QFT}

So what is the full quantum Fourier transform? It is the transform
which takes the amplitudes of a $N$ dimensional state and computes the
Fourier transform on these amplitudes (which are then the new
amplitudes in the computational basis.) In other words, the QFT enacts
the transform

\begin{equation*}
\sum_{x=0}^{N-1} a_{x}|x\rangle \rightarrow \sum_{x=0}^{N-1} \tilde{a}_{x}|x\rangle=\sum_{x=0}^{N-1} \frac{1}{\sqrt{N}} \sum_{y=0}^{N-1} \omega_{N}^{-x y} a_{y}|x\rangle 
\end{equation*}
\end{frame}

\begin{frame}[plain,fragile]
\frametitle{Explicit transform}

It is easy to see that this implies that the QFT performs the following transform on basis states:

\begin{equation*}
|x\rangle \rightarrow \frac{1}{\sqrt{N}} \sum_{y=0}^{N-1} \omega_{N}^{-x y}|y\rangle
\end{equation*}

Thus the QFT is given by the matrix

\begin{equation*}
U_{Q F T}=\frac{1}{\sqrt{N}} \sum_{x=0}^{N-1} \sum_{y=0}^{N-1} \omega_{N}^{-y x}|y\rangle\langle x| 
\end{equation*}
\end{frame}

\begin{frame}[plain,fragile]
\frametitle{Unitarity}

The last  matrix is unitary. Let's check this:

\begin{align*}
U_{Q F T} U_{Q F T}^{\dagger} & =\frac{1}{N} \sum_{x=0}^{N-1} \sum_{y=0}^{N-1} \omega_{N}^{y x}|x\rangle\left\langle y\left|\sum_{x^{\prime}=0}^{N-1} \sum_{y^{\prime}=0}^{N-1} \omega_{N}^{-y^{\prime} x^{\prime}}\right| y^{\prime}\right\rangle\left\langle x^{\prime}\right| \\
& =\frac{1}{N} \sum_{x, y, x^{\prime}, y^{\prime}=0}^{N-1} \omega_{N}^{y x-y^{\prime} x^{\prime}} \delta_{y, y^{\prime}}|x\rangle\left\langle x^{\prime}\left|=\frac{1}{N} \sum_{x, y, x^{\prime}=0}^{N-1} \omega_{N}^{y\left(x-x^{\prime}\right)}\right| x\right\rangle\left\langle x^{\prime}\right| \\
& =\sum_{x, x^{\prime}=0}^{N-1} \delta_{x, x^{\prime}}|x\rangle\left\langle x^{\prime}\right|=I 
\end{align*}
\end{frame}

\begin{frame}[plain,fragile]
\frametitle{Importance of QFT}

The QFT is a very important transform in quantum computing. It can be
used for all sorts of cool tasks, including, as we shall see in Shor's
algorithm. But before we can use it for quantum computing tasks, we
should try to see if we can efficiently implement the QFT with a
quantum circuit. Indeed we can and the reason we can is intimately
related to the fast Fourier transform.
\end{frame}

\begin{frame}[plain,fragile]
\frametitle{Circuit QFT}

Let's derive a circuit for the QFT when $N=2^{n}$. The QFT performs the transform

\begin{equation*}
|x\rangle \rightarrow \frac{1}{\sqrt{2^{n}}} \sum_{y=0}^{2^{n}-1} \omega_{N}^{-x y}|y\rangle 
\end{equation*}

Then we can expand out this sum

\begin{equation*}
|x\rangle \rightarrow \frac{1}{\sqrt{2^{n}}} \sum_{y_{1}, y_{2}, \ldots, y_{n} \in\{0,1\}} \omega_{N}^{-x \sum_{k=1}^{n} 2^{n-k} y_{k}}\left|y_{1}, y_{2}, \ldots, y_{n}\right\rangle 
\end{equation*}
\end{frame}

\begin{frame}[plain,fragile]
\frametitle{Expanding the exponential}

Expanding the exponential of a sum to a product of exponentials and
collecting these terms in from the appropriate terms we can express
this as

\begin{equation*}
|x\rangle \rightarrow \frac{1}{\sqrt{2^{n}}} \sum_{y_{1}, y_{2}, \ldots, y_{n} \in\{0,1\}} \bigotimes_{k=1}^{n} \omega_{N}^{-x 2^{n-k} y_{k}}\left|y_{k}\right\rangle 
\end{equation*}
\end{frame}

\begin{frame}[plain,fragile]
\frametitle{Rearranging}

We can rearrange the sum and products as

\begin{equation*}
|x\rangle \rightarrow \frac{1}{\sqrt{2^{n}}} \bigotimes_{k=1}^{n}\left[\sum_{y_{k} \in\{0,1\}} \omega_{N}^{-x 2^{n-k} y_{k}}\left|y_{k}\right\rangle\right] 
\end{equation*}

Expanding this sum yields

\begin{equation*}
|x\rangle \rightarrow \frac{1}{\sqrt{2^{n}}} \bigotimes_{k=1}^{n}\left[|0\rangle+\omega_{N}^{-x 2^{n-k}}|1\rangle\right]
\end{equation*}
\end{frame}

\begin{frame}[plain,fragile]
\frametitle{Notation for binary fraction}

But now notice that $\omega_{N}^{-x 2^{n-k}}$ is not dependent on the
higer order bits of $x$. It is convenient to adopt the following
expression for a binary fraction:

\begin{equation*}
0 . x_{l} x_{l+1} \ldots x_{n}=\frac{x_{l}}{2}+\frac{x_{l+1}}{4}+\cdots+\frac{x_{n}}{2^{n-l+1}}
\end{equation*}
\end{frame}

\begin{frame}[plain,fragile]
\frametitle{More notations}

Then we can see that
\begin{equation*}
|x\rangle \rightarrow \frac{1}{\sqrt{2^{n}}}\left[|0\rangle+e^{-2 \pi i 0 . x_{n}}|1\rangle\right] \otimes\left[|0\rangle+e^{-2 \pi i 0 . x_{n-1} x_{n}}|1\rangle\right] \otimes \cdots \otimes\left[|0\rangle+e^{-2 \pi i 0 . x_{1} x_{2} \cdots x_{n}}|1\rangle\right] 
\end{equation*}

This is a very useful form of the QFT for $N=2^{n}$. Why? Because we
see that only the last qubit depends on the the values of all the
other input qubits and each further bit depends less and less on the
input qubits. Further we note that $e^{-2 \pi i 0 . a}$ is either +1
or -1 , which reminds us of the Hadamard transform.
\end{frame}

\begin{frame}[plain,fragile]
\frametitle{Deriving a circuit}

So how do we use this to derive a circuit for the QFT over $N=2^{n}$ ?

Take the first qubit of $\left|x_{1}, \ldots, x_{n}\right\rangle$ and
apply a Hadamard transform. This produces the transform

\begin{equation*}
|x\rangle \rightarrow \frac{1}{\sqrt{2}}\left[|0\rangle+e^{-2 \pi i 0 \cdot x_{1}}|1\rangle\right] \otimes\left|x_{2}, x_{3}, \ldots, x_{n}\right\rangle 
\end{equation*}
\end{frame}

\begin{frame}[plain,fragile]
\frametitle{Rotation gate}

Now define the rotation gate

\[
R_{k}=\left[\begin{array}{cc}
1 & 0  \tag{30}\\
0 & \exp \left(\frac{-2 \pi i}{2^{k}}\right)
\end{array}\right]
\]
If we now apply controlled $R_{2}, R_{3}$, etc. gates controlled on
the appropriate bits this enacts the transform

\begin{equation*}
|x\rangle \rightarrow \frac{1}{\sqrt{2}}\left[|0\rangle+e^{-2 \pi i 0 . x_{1} x_{2} \ldots x_{n}}|1\rangle\right] \otimes\left|x_{2}, x_{3}, \ldots, x_{n}\right\rangle 
\end{equation*}
\end{frame}

\begin{frame}[plain,fragile]
\frametitle{Proceeding}

Thus we have reproduced the last term in the QFTed state. Of course
now we can proceed to the second qubit, perform a Hadamard, and the
appropriate controlled $R_{k}$ gates and get the second to last
qubit. Thus when we are finished we will have the transform

\begin{equation*}
|x\rangle \rightarrow \frac{1}{\sqrt{2^{n}}}\left[|0\rangle+e^{-2 \pi i 0 . x_{1} x_{2} \cdots x_{n}}|1\rangle\right] \otimes\left[|0\rangle+e^{-2 \pi i 0 . x_{1} x_{2} \cdots x_{n-1}}|1\rangle\right] \otimes \cdots \otimes\left[|0\rangle+e^{-2 \pi i 0 . x_{n}}|1\rangle\right]
\end{equation*}

Reversing the order of these qubits will then produce the QFT!

% The circuit we have constructed on $n$ qubits is

% \includegraphics[max width=\textwidth]{2024_03_18_c1d427fa6b85efa365f1g-5}
\end{frame}

\begin{frame}[plain,fragile]
\frametitle{Two-qubit case}

For a two-qubit system $n=2$, the QFT transformation is:
\[
\text{QFT} \vert x_0 x_1\rangle = \frac{1}{2} \sum_{y=0}^{3} e^{2\pi i x y / 4} \vert y\rangle, 
\]
where $x = x_0 x_1$ (binary representation).
\end{frame}

\begin{frame}[plain,fragile]
\frametitle{QFT matrix}

The QFT on two qubits has the matrix:
\[
   U_{\text{QFT}} = \frac{1}{2}
   \begin{bmatrix}
       1 & 1 & 1 & 1 \\
       1 & i & -1 & -i \\
       1 & -1 & 1 & -1 \\
       1 & -i & -1 & i
   \end{bmatrix}
\]
Exercise: apply this to the $\vert 01\rangle$ (binary 1) and $\vert 10\rangle$ (binary 2) stats.
Hint: see Hundt section 6.2 and circuit figures as well.
\end{frame}

\begin{frame}[plain,fragile]
\frametitle{Four-qubit system}

The basis states are
\[
|j_1j_2j_3j_4 \rangle
\]
where \(j_k\) is either \(0\) or \(1\). We have

\begin{align*}
0.j_3 &= \frac{j_3}{2} \\
0.j_2j_3 &= \frac{j_2}{2} + \frac{j_3}{4} \\
0.j_1j_2j_3 &= \frac{j_1}{2} + \frac{j_2}{4} + \frac{j_3}{8} \\
0.j_0j_1j_2j_3 &= \frac{j_0}{2} + \frac{j_1}{4} + \frac{j_2}{8} + \frac{j_3}{16} \\
\end{align*}
\end{frame}

\begin{frame}[plain,fragile]
\frametitle{QFT acts as follows}

The quantum Fourier transform acts as follows:

\begin{align*}
|j_1j_2j_3j_4 \rangle \mapsto \frac{1}{\sqrt{2^{4/2}}}
\left(|0\rangle + e^{2 \pi i 0.j4}|1\rangle \right) \otimes 
\left(|0\rangle + e^{2 \pi i 0.j_3j4}|1\rangle \right) \otimes 
\left(|0\rangle + e^{2 \pi i 0.j_2j_3j4}|1\rangle \right) \otimes 
\left(|0\rangle + e^{2 \pi i 0.j_1j_2j_3j4}|1\rangle \right) 
\end{align*}
\end{frame}

\begin{frame}[plain,fragile]
\frametitle{Quantum circuits}

To compose a quantum circuit that calculates the quantum Fourier
transform we use the operators

\begin{align*}
R_k = 
\begin{bmatrix}
1 & 0 \\
0 & e^{2\pi i/2^k}
\end{bmatrix}.
\end{align*}
\end{frame}

\begin{frame}[plain,fragile]
\frametitle{Rotation gates}

In this example, the \(R_k\) gates are:
\begin{align*}
R_1 = \begin{bmatrix}
1 & 0 \\
0 & e^{2 \pi i /2^0}
\end{bmatrix} = 
\begin{bmatrix}
1 & 0 \\
0 & 1
\end{bmatrix}, \quad 
R_2 = \begin{bmatrix}
1 & 0 \\
0 & e^{2 \pi i /2^2}
\end{bmatrix}, \quad 
R_3 = \begin{bmatrix}
1 & 0 \\
0 & e^{2 \pi i /2^3}
\end{bmatrix}, \quad 
R_4 = \begin{bmatrix}
1 & 0 \\
0 & e^{2 \pi i /2^4}
\end{bmatrix}.
\end{align*}
\end{frame}

\begin{frame}[plain,fragile]
\frametitle{Using the Hadamard gate}

The Hadamard
gate on a single qubit creates an equal superposition of its basis
states, assuming it is not already in a superposition, such that

\[
H\vert 0 \rangle = \frac{1}{\sqrt{2}} \left(\vert 0 \rangle + \vert 1\rangle\right), \quad H\vert 1\rangle = \frac{1}{\sqrt{2}} \left(\vert 0 \rangle - \vert 1\rangle\right)
\]
The $R_k$ gate simply adds a phase if the qubit it acts on is in the state $\vert 1\rangle$
\[
    R_k\vert 0 \rangle = \vert 0 \rangle, \quad R_k\vert 1\rangle = e^{2\pi i/2^{k}}\vert 1\rangle
\]

Since all this gates are unitary, the quantum Fourier transfrom is also unitary.
\end{frame}

\begin{frame}[plain,fragile]
\frametitle{Algorithm}

Assume we have a quantum register of $n$ qubits in the state $\vert j_1 j_2 \dots j_n\rangle$.
Applying the Hadamard gate to the first qubit
produces the state

\[
H\vert j_1 j_2 \dots j_n\rangle = \frac{\left(\vert 0 \rangle + e^{2\pi i 0.j_1}\vert 1\rangle\right)}{2^{1/2}} \vert j_2 \dots j_n\rangle.
\]
\end{frame}

\begin{frame}[plain,fragile]
\frametitle{Binary fraction}

Here we have made use of the binary fraction to represent the action of the Hadamard gate 
\[
\exp{2\pi i 0.j_1} = -1,
\]
if $j_1 = 1$ and $+1$ if $j_1 = 0$.
\end{frame}

\begin{frame}[plain,fragile]
\frametitle{Controlled rotation gate}

Furthermore we can apply the controlled-$R_k$ gate, with all the other qubits $j_k$ for $k>1$ as control qubits to produce the state

\[
\frac{\left(\vert 0 \rangle + e^{2\pi i 0.j_1j_2\dots j_n}\vert 1\rangle\right)}{2^{1/2}} \vert j_2 \dots j_n\rangle
\]

Next we do the same procedure on qubit $2$ producing the state

\[
\frac{\left(\vert 0 \rangle + e^{2\pi i 0.j_1j_2\dots j_n}\vert 1\rangle\right)\left(\vert 0 \rangle + e^{2\pi i 0.j_2\dots j_n}\vert 1\rangle\right)}{2^{2/2}} \vert j_2 \dots j_n\rangle
\]
\end{frame}

\begin{frame}[plain,fragile]
\frametitle{Applying to all qubits}

Doing this for all $n$ qubits yields state

\[
\frac{\left(\vert 0 \rangle + e^{2\pi i 0.j_1j_2\dots j_n}\vert 1\rangle\right)\left(\vert 0 \rangle + e^{2\pi i 0.j_2\dots j_n}\vert 1\rangle\right)\dots \left(\vert 0 \rangle + e^{2\pi i 0.j_n}\vert 1\rangle\right)}{2^{n/2}} \vert j_2 \dots j_n\rangle
\]
At the end we use swap gates to reverse the order of the qubits

\[
\frac{\left(\vert 0 \rangle + e^{2\pi i 0.j_n}\vert 1\rangle\right)\left(\vert 0 \rangle + e^{2\pi i 0.j_{n-1}j_n}\vert 1\rangle\right)\dots\left(\vert 0 \rangle + e^{2\pi i 0.j_1j_2\dots j_n}\vert 1\rangle\right) }{2^{n/2}} \vert j_2 \dots j_n\rangle
\]
This is just the product representation from earlier, obviously our desired output.
\end{frame}




%------------------------------------------------

\begin{frame}{Interpretation}

Each qubit accumulates phases from less significant bits.

Controlled rotations naturally appear.

\end{frame}



\begin{frame}{General QFT Circuit}

\begin{center}
\begin{quantikz}
\lstick{q0}&\gate{H}&\ctrl{1}&\ctrl{2}&\swap{2}&\qw\\
\lstick{q1}&\qw&\gate{R2}&\ctrl{1}&\qw&\qw\\
\lstick{q2}&\qw&\qw&\gate{R3}&\gate{H}&\swap{-2}
\end{quantikz}
\end{center}

\end{frame}


%------------------------------------------------

\begin{frame}{Two-Qubit QFT}

\begin{center}
\begin{quantikz}
\lstick{}&\gate{H}&\ctrl{1}&\swap{1}\\
\lstick{}&\gate{R2}&\gate{H}&\swap{-1}
\end{quantikz}
\end{center}

\end{frame}

%------------------------------------------------

\begin{frame}{Three-Qubit QFT}

\begin{center}
\begin{quantikz}
\lstick{}&\gate{H}&\ctrl{1}&\ctrl{2}&\swap{2}\\
\lstick{}&\gate{R2}&\gate{H}&\ctrl{1}&\qw\\
\lstick{}&\gate{R3}&\gate{R2}&\gate{H}&\swap{-2}
\end{quantikz}
\end{center}

\end{frame}

%------------------------------------------------

\begin{frame}{Gate Complexity}

Number of gates

\[
\frac{n(n+1)}{2}
\]

Scaling

\[
O(n^2)
\]

\end{frame}

%------------------------------------------------

\begin{frame}{Approximate QFT}

Small rotations contribute little.

Truncate rotations

\[
R_k \quad k>m
\]


Approximate QFT

\[
O(n\log n)
\]

\end{frame}

\begin{frame}[plain,fragile]
\frametitle{Plans for next week, March 23-27}

\begin{enumerate}
\item From QFTs to the quantum phase estimation algorithm and eigenvalues of for example the Lipkin  Hamiltonian

\item Discussions of variants for project 2
\end{enumerate}

\noindent
\end{frame}

\end{document}



